\documentclass[11pt]{article}
\usepackage[margin=1in]{geometry}
\usepackage[T1]{fontenc}
\usepackage[utf8]{inputenc}
\usepackage{amsmath,amssymb,amsthm}
\usepackage{enumitem}

\title{ICPC World Finals 2021\\J. Splitstream}
\author{}
\date{}

\begin{document}
\maketitle

\section*{Problem Summary}

The network starts with the sequence
\[
    1,2,3,\dots,m.
\]
Each node either:
\begin{itemize}[leftmargin=*]
    \item splits one sequence into odd and even positions, or
    \item merges two sequences by alternating their elements.
\end{itemize}

For each query $(x,k)$, we must report the $k$th value on output wire $x$, or \texttt{none} if that
output is shorter than $k$.

\section*{Main Observation}

We never need to materialize any sequence.

If we know the length of every wire, then for a query $(x,k)$ we can walk \emph{backwards} through the
network until we reach the original input wire $1$.

\section*{Lengths of Wires}

Let $\ell(w)$ be the length of wire $w$.

\paragraph{Split node.}
If a split node reads input $x$ and writes outputs $y$ and $z$, then:
\[
    \ell(y)=\left\lceil\frac{\ell(x)}{2}\right\rceil,
    \qquad
    \ell(z)=\left\lfloor\frac{\ell(x)}{2}\right\rfloor.
\]

\paragraph{Merge node.}
If a merge node reads inputs $x$ and $y$ and writes output $z$, then:
\[
    \ell(z)=\ell(x)+\ell(y).
\]

These values can be computed with memoized DFS because the network is acyclic.

\section*{Tracing One Query Backwards}

\subsection*{Split Output}

Suppose a split node takes input sequence
\[
    a_1,a_2,a_3,\dots
\]
Then:
\begin{itemize}[leftmargin=*]
    \item output 1 is
    \[
        a_1,a_3,a_5,\dots
    \]
    \item output 2 is
    \[
        a_2,a_4,a_6,\dots
    \]
\end{itemize}

So:
\begin{itemize}[leftmargin=*]
    \item the $k$th element of output 1 is the $(2k-1)$th element of the input;
    \item the $k$th element of output 2 is the $(2k)$th element of the input.
\end{itemize}

\subsection*{Merge Output}

Suppose a merge node takes sequences
\[
    a_1,a_2,\dots,a_p
    \qquad\text{and}\qquad
    b_1,b_2,\dots,b_q.
\]
Its output is
\[
    a_1,b_1,a_2,b_2,\dots
\]
until one sequence ends, and then the remaining elements of the longer sequence continue unchanged.

Therefore:
\begin{itemize}[leftmargin=*]
    \item while $k \le 2\min(p,q)$:
    \begin{itemize}[leftmargin=*]
        \item odd $k$ comes from the first input at index $(k+1)/2$;
        \item even $k$ comes from the second input at index $k/2$.
    \end{itemize}
    \item after that, if $p>q$, the remaining elements come from the first input at index $k-q$;
    \item if $q>p$, they come from the second input at index $k-p$.
\end{itemize}

\section*{Algorithm}

\begin{enumerate}[leftmargin=*]
    \item Parse all nodes and record, for each output wire, which node created it.
    \item Compute wire lengths lazily with memoized DFS.
    \item For each query $(x,k)$:
    \begin{itemize}[leftmargin=*]
        \item if $k > \ell(x)$, print \texttt{none};
        \item otherwise repeatedly replace $(x,k)$ by the corresponding predecessor wire and predecessor
        index, using the rules above, until $x=1$.
    \end{itemize}
    \item On wire $1$, the sequence is simply $1,2,\dots,m$, so the answer is $k$ itself.
\end{enumerate}

\section*{Correctness Proof}

We prove that the algorithm answers every query correctly.

\paragraph{Lemma 1.}
For every wire, the memoized formulas compute its correct length.

\paragraph{Proof.}
The formulas follow directly from the definitions of split and merge nodes:
\begin{itemize}[leftmargin=*]
    \item a split sends odd-positioned elements to one output and even-positioned elements to the other;
    \item a merge outputs every element from both inputs exactly once.
\end{itemize}
Since the network is acyclic, recursively applying these formulas reaches the base wire $1$ and is
well-defined. \qed

\paragraph{Lemma 2.}
For a split node, the backward index transformation used by the algorithm is correct.

\paragraph{Proof.}
By definition of the split operation, output 1 contains exactly the odd-indexed elements of the input, in
the same order. So its $k$th element is input element $2k-1$. Likewise output 2 contains exactly the
even-indexed input elements, so its $k$th element is input element $2k$. \qed

\paragraph{Lemma 3.}
For a merge node, the backward index transformation used by the algorithm is correct.

\paragraph{Proof.}
The merge node alternates elements from the two inputs as long as both still have elements. Therefore the
first $2\min(p,q)$ output positions correspond exactly to alternating positions from the two inputs. After
the shorter input is exhausted, the output continues with the remaining suffix of the longer input without
any further mixing. The backward formulas are exactly these cases written explicitly. \qed

\paragraph{Theorem.}
For every query $(x,k)$, the algorithm outputs the correct $k$th element of wire $x$, or \texttt{none} if
that element does not exist.

\paragraph{Proof.}
If $k > \ell(x)$, then wire $x$ has fewer than $k$ elements, so \texttt{none} is correct by Lemma 1.

Otherwise, by repeatedly applying Lemmas 2 and 3, the algorithm transforms the query to an equivalent
query on the predecessor wire of the current node. Because the network is acyclic, this process
eventually reaches wire $1$. On wire $1$, the $k$th element is exactly $k$. Thus the returned value is
the unique value that maps through the network to the original query position. \qed

\section*{Complexity Analysis}

Let $N$ be the number of nodes and $Q$ the number of queries.

Each wire length is computed at most once, so the total preprocessing is $O(N)$. Each query walks upward
through at most one path in the DAG, so one query costs $O(N)$ in the worst case, which is easily fast
enough for the given limits. Memory usage is $O(N)$.

\section*{Implementation Notes}

\begin{itemize}[leftmargin=*]
    \item Wire lengths fit in 64-bit integers because no wire can contain more than the original $m$
    values.
    \item The answer on wire $1$ is just the queried index itself.
\end{itemize}

\end{document}
